\section{Conclusion}
\label{sec:conclusion}

We presented a fully automated multi-panel story image generation system centered on a \textbf{subject-count hybrid router}.
Single-character narratives use structured LLM planning, run-local anchor banks, IP-Adapter conditioning on SDXL-Turbo, and scene-consistency prompt assembly, with identity strength $\lambda=0.3$ as a practical default for Story-14-style results.
Multi-character narratives route to native StoryDiffusion to avoid single-anchor face contamination while preserving separate identity banks for each tagged entity.

We also implemented an exploratory PixArt-$\alpha$ DiT backend with DreamBooth LoRA and cross-scene attention blending, motivated by feedback suggesting transformer-based architectures. The engineering migration was completed---including custom attention wrappers, LoRA training on augmented generated images, and T5 trigger-token binding---but qualitative results did not meet the SDXL + IP-Adapter baseline. Through detailed failure analysis (Section~\ref{sec:dit_discussion}), we identified specific root causes: DreamBooth trained on generated images learns generator-specific artifacts rather than subject identity; augmentation-based training lacks pose and context diversity; long inference prompts dilute the trigger token's cross-attention weight; and cross-scene blending partially overwrites LoRA identity features in middle-to-deep transformer layers. We document these findings as a rigorous negative result rather than omitting the implementation effort, and identify concrete directions for future DiT-based approaches: PixArt-compatible IP-Adapter weights, text-encoder LoRA with dedicated trigger-token embedding optimization, and late-step-only cross-scene attention.

The central lesson is architectural: separate \emph{text understanding}, \emph{safe identity mechanisms}, and \emph{backend choice by subject count}, with logged deterministic decisions, yields more reliable story panels than any single generator applied uniformly to all stories. The DiT exploration further teaches that the ecosystem maturity around a backbone---adapter availability, distillation quality, community tooling---can outweigh the theoretical advantages of a newer architecture for a fixed-scope course project.
